% !TeX program = xelatex
\documentclass[a4paper,12pt]{article}

% ===== 1. KHO GIAY, LE TRANG, FONT =====
\usepackage[left=3cm,right=2cm,top=2cm,bottom=2cm,footskip=1cm]{geometry}
\usepackage{fontspec}

\IfFontExistsTF{Times New Roman}{
  \setmainfont{Times New Roman}
}{
  \IfFileExists{fonts/times.ttf}{
    \setmainfont{times.ttf}[
      Path=fonts/,
      BoldFont=timesbd.ttf,
      ItalicFont=timesi.ttf,
      BoldItalicFont=timesbi.ttf]
  }{
    \IfFontExistsTF{TeX Gyre Termes}{
      \setmainfont{TeX Gyre Termes}
      \PackageWarningNoLine{MauBaoCao}{Times New Roman unavailable; using TeX Gyre Termes}
    }{
      \setmainfont{DejaVu Serif}
      \PackageWarningNoLine{MauBaoCao}{Times New Roman unavailable; using DejaVu Serif}
    }
  }
}
\usepackage{indentfirst}
\usepackage{microtype}
\usepackage{graphicx}
\usepackage{tikz}
\usepackage{array,tabularx,booktabs,longtable}
\usepackage{enumitem}
\usepackage{titlesec}
\usepackage{tocloft}
\usepackage{caption}
\usepackage{xurl}
\usepackage[unicode,hidelinks]{hyperref}
\urlstyle{same}

\renewcommand{\normalsize}{\fontsize{13}{19.5}\selectfont}
\AtBeginDocument{\normalsize}
\setlength{\parindent}{0.5cm}
\setlength{\parskip}{4pt}
\setlength{\emergencystretch}{2em}
\raggedbottom
\setlist{topsep=4pt,itemsep=2pt,parsep=0pt,leftmargin=*}

% ===== 2. BA CAP TIEU DE: 1 / 1.1 / 1.1.1 =====
\setcounter{secnumdepth}{3}
\setcounter{tocdepth}{3}
\renewcommand{\thesection}{\arabic{section}}
\renewcommand{\thesubsection}{\thesection.\arabic{subsection}}
\renewcommand{\thesubsubsection}{\thesubsection.\arabic{subsubsection}}
\titleformat{\section}[hang]
  {\normalfont\fontsize{16}{20}\selectfont\bfseries}
  {\thesection.}{0.6em}{}
\titleformat{\subsection}[hang]
  {\normalfont\fontsize{14}{18}\selectfont\bfseries}
  {\thesubsection}{0.6em}{}
\titleformat{\subsubsection}[hang]
  {\normalfont\fontsize{13}{17}\selectfont\bfseries}
  {\thesubsubsection}{0.6em}{}
\titlespacing*{\section}{0pt}{16pt}{8pt}
\titlespacing*{\subsection}{0pt}{12pt}{6pt}
\titlespacing*{\subsubsection}{0pt}{8pt}{4pt}

% ===== 3. MUC LUC VA SO TRANG =====
\renewcommand{\contentsname}{MỤC LỤC}
\renewcommand{\refname}{Tài liệu tham khảo}
\renewcommand{\cfttoctitlefont}{\fontsize{16}{20}\selectfont\bfseries}
\setlength{\cftbeforetoctitleskip}{0pt}
\setlength{\cftaftertoctitleskip}{18pt}
\renewcommand{\cftsecfont}{\normalfont\bfseries}
\renewcommand{\cftsecpagefont}{\normalfont\bfseries}
\renewcommand{\cftsubsecfont}{\normalfont}
\renewcommand{\cftsubsubsecfont}{\normalfont}
\renewcommand{\cftsecleader}{\cftdotfill{\cftdotsep}}
\cftsetindents{section}{0em}{2em}
\cftsetindents{subsection}{2em}{3em}
\cftsetindents{subsubsection}{5em}{4em}
\setlength{\cftbeforesecskip}{6pt}
\makeatletter
\def\ps@plain{%
  \let\@oddhead\@empty\let\@evenhead\@empty
  \def\@oddfoot{\hfil{\normalfont\fontsize{11}{13}\selectfont\thepage}\hfil}%
  \let\@evenfoot\@oddfoot}
\makeatother
\pagestyle{plain}

% ===== 4. BANG, HINH, CHU THICH =====
\renewcommand{\tablename}{Bảng}
\renewcommand{\figurename}{Hình}
\DeclareCaptionFont{reportcaption}{\fontsize{12}{15}\selectfont}
\captionsetup{font=reportcaption,labelfont=bf,labelsep=period,
  justification=centering,skip=6pt}
\newcolumntype{Y}{>{\raggedright\arraybackslash}X}
\newcommand{\bangfont}{\fontsize{12}{15}\selectfont}

% ===== 5. THONG TIN CAN THAY TREN BIA =====
\newcommand{\DonViChuQuan}{ĐẠI HỌC QUỐC GIA HÀ NỘI}
\newcommand{\TenTruong}{TRƯỜNG ĐẠI HỌC CÔNG NGHỆ}
\newcommand{\LoaiBaoCao}{PROGRESSION REPORT 1}
\newcommand{\TenMonHoc}{ĐIỆN TOÁN ĐÁM MÂY}
\newcommand{\TenDeTai}{SatiBot - Quản lý và tương tác với các tài nguyên AWS bằng ngôn ngữ tự nhiên}
\newcommand{\GiangVien}{Hoàng Xuân Tùng}
\newcommand{\SoNhom}{10}
\newcommand{\LopMonHoc}{INT3319 1}
\newcommand{\DiaDiemThoiGian}{Hà Nội, Tháng 10/2026}
\hypersetup{pdftitle={PROGRESSION REPORT 1 - BÁO CÁO TIẾN ĐỘ TUẦN 1},pdfauthor={Nhóm thực hiện}}

\begin{document}

% ===== 6. TRANG BIA =====
\begin{titlepage}
\thispagestyle{empty}
\begin{tikzpicture}[remember picture,overlay]
  \draw[line width=2pt]
    ([xshift=3cm,yshift=-2cm]current page.north west)
    rectangle ([xshift=-2cm,yshift=2cm]current page.south east);
  \draw[line width=0.5pt]
    ([xshift=3.12cm,yshift=-2.12cm]current page.north west)
    rectangle ([xshift=-2.12cm,yshift=2.12cm]current page.south east);
\end{tikzpicture}\par
\centering
\setlength{\parindent}{0pt}
\setlength{\parskip}{0pt}
\vspace*{0.8cm}
{\fontsize{14}{18}\selectfont\bfseries\DonViChuQuan\par}
\vspace{0.15cm}
{\fontsize{15}{20}\selectfont\bfseries\TenTruong\par}
\vspace{0.25cm}
\rule{5cm}{0.4pt}\par
\vspace{0.5cm}

\IfFileExists{logo_uet.png}{
  \includegraphics[width=3.8cm,height=3.8cm,keepaspectratio]{logo_uet.png}
}{
  \IfFileExists{logo.png}{
    \includegraphics[width=3.8cm,height=3.8cm,keepaspectratio]{logo.png}
  }{
    \fbox{\parbox[c][3.4cm][c]{3.4cm}{\centering
      \fontsize{11}{14}\selectfont LOGO TRƯỜNG\\(logo.png)}}
  }
}\par
\vspace{0.7cm}
{\fontsize{22}{27}\selectfont\bfseries\LoaiBaoCao\par}
\vspace{0.35cm}
\begin{minipage}{0.9\textwidth}\centering
  {\fontsize{15}{20}\selectfont\bfseries MÔN HỌC: \TenMonHoc\par}
\end{minipage}\par
\vspace{0.6cm}
{\fontsize{13}{17}\selectfont\bfseries ĐỀ TÀI\par}
\vspace{0.12cm}
\begin{minipage}{0.9\textwidth}\centering
  {\fontsize{16}{21}\selectfont\bfseries\TenDeTai\par}
\end{minipage}\par
\vspace{0.65cm}
{\fontsize{13}{19.5}\selectfont
\renewcommand{\arraystretch}{1.1}
\begin{tabular}{@{}>{\raggedright\arraybackslash\bfseries}p{4.8cm}>{\raggedright\arraybackslash}p{8cm}@{}}
  Giảng viên hướng dẫn: & \GiangVien\\
  Nhóm thực hiện: & \SoNhom\\
  Sinh viên thực hiện: & Cao Anh Minh \textbf{(Trưởng nhóm)}\\
                      & Trịnh Thiên Quang\\
                      & Phạm Đức Hùng\\
                      & Nguyễn Đức Huy\\
  Lớp môn học: &  \LopMonHoc
\end{tabular}\par}
\vfill
{\fontsize{13}{17}\selectfont\bfseries\DiaDiemThoiGian\par}
\vspace*{0.5cm}
\end{titlepage}

% ===== 7. MUC LUC TU DONG =====
\pagenumbering{roman}
\tableofcontents
\clearpage
\pagenumbering{arabic}

% ===== 8. NOI DUNG BAO CAO =====

\section{TỔNG QUAN HỆ THỐNG}

\subsection{Ý tưởng và mục tiêu}

\subsubsection{Bài toán và giải pháp}
Trong quá trình quản trị và vận hành hạ tầng trên nền tảng điện toán đám mây Amazon Web Services (AWS), việc thao tác qua giao diện đồ họa (AWS Management Console) hoặc giao diện dòng lệnh (AWS CLI) đòi hỏi người dùng phải ghi nhớ nhiều thao tác phức tạp và cú pháp kịch bản chính xác. Để giải quyết rào cản này, nhóm đề xuất xây dựng một \textbf{ứng dụng web ChatOps} hỗ trợ người dùng quản lý và tương tác với các tài nguyên AWS bằng ngôn ngữ tự nhiên, tối ưu hóa cho tiếng Việt.

\subsubsection{Giá trị dự kiến}
Sản phẩm hướng tới việc đơn giản hóa quy trình quản trị điện toán đám mây, nâng cao trải nghiệm vận hành và giảm thiểu sai sót do thao tác thủ công.

\subsection{Mô tả hệ thống}

\subsubsection{Hình thức sản phẩm}
Sản phẩm được định hình là một ứng dụng web.

\subsubsection{Quy trình trải nghiệm người dùng}
Trải nghiệm sử dụng hệ thống của người dùng diễn ra qua các bước sau:
\begin{enumerate}
  \item \textbf{Đăng nhập:} Người dùng truy cập website qua trình duyệt và đăng nhập vào hệ thống bằng danh tính được cấp.
  \item \textbf{Nhập yêu cầu:} Người dùng nhập câu lệnh yêu cầu bằng ngôn ngữ tự nhiên (ưu tiên tiếng Việt) trên giao diện trò chuyện.
  \item \textbf{Bổ sung thông tin (nếu cần):} Hệ thống phân tích ý định và các tham số. Nếu phát hiện thiếu thông tin hoặc chưa xác định rõ tài nguyên đích, hệ thống sẽ hỏi lại để người dùng làm rõ.
  \item \textbf{Xác nhận kế hoạch:} Hệ thống trình bày kế hoạch thao tác dự kiến và yêu cầu người dùng xác nhận đối với các tác vụ cần thiết.
  \item \textbf{Thực thi và nhận kết quả:} Sau khi được xác nhận, backend kiểm tra quyền, tham số và thực hiện thao tác tương ứng trên AWS. Kết quả được phản hồi trực tiếp trên giao diện, đồng thời lưu lại lịch sử và trạng thái tác vụ. Trong trường hợp thao tác thất bại hoặc chỉ hoàn thành một phần, hệ thống sẽ thông báo chính xác tình trạng đó cho người dùng.
\end{enumerate}

\subsection{Phạm vi}

\subsubsection{Các dịch vụ AWS được quản lý}
Hệ thống tập trung xử lý các tác vụ quản trị phục vụ người dùng đối với 04 dịch vụ AWS nòng cốt:
\begin{itemize}
  \item \textbf{IAM} (Identity and Access Management)
  \item \textbf{S3} (Simple Storage Service)
  \item \textbf{EC2} (Elastic Compute Cloud)
  \item \textbf{Lambda} (AWS Lambda)
\end{itemize}

\subsubsection{Định hướng mở rộng}
Sau khi hoàn thiện cơ chế quản lý ổn định cho 4 dịch vụ trên, hệ thống có thể được nghiên cứu mở rộng thêm các dịch vụ AWS khác trong các giai đoạn phát triển tiếp theo.


\section{NGƯỜI DÙNG VÀ KỊCH BẢN SỬ DỤNG}

\subsection{Đối tượng sử dụng}
Hệ thống được thiết kế hướng tới các nhóm người dùng mục tiêu sau:
\begin{itemize}
  \item \textbf{Nhân viên kỹ thuật và lập trình viên:} Những người có nhu cầu thiết lập, cấu hình hoặc vận hành các tài nguyên AWS trong quá trình phát triển sản phẩm.
  \item \textbf{Sinh viên:} Luyện tập và thực hành các thao tác quản trị môn Điện toán đám mây.
  \item \textbf{Chủ doanh nghiệp:} Có nhu cầu theo dõi, quản lý tài nguyên cloud của đơn vị mình thông qua các chức năng phù hợp với phạm vi quyền hạn được cấp.
\end{itemize}

\subsection{Các kịch bản sử dụng tập trung}

\subsubsection{Nhóm kịch bản quản lý S3}
\begin{enumerate}
  \item Tạo bucket riêng tư.
  \item Tra cứu bucket, danh sách file và dung lượng.
  \item Chia sẻ file bằng liên kết tải có thời hạn.
  \item Kiểm tra và hạn chế truy cập công khai.
  \item Xóa bucket và dữ liệu có kiểm soát.
\end{enumerate}

\subsubsection{Nhóm kịch bản quản lý IAM}
\begin{enumerate}
  \item Tra cứu role và policy.
  \item Tạo role hoặc policy.
  \item Cấp hoặc điều chỉnh quyền truy cập.
  \item Thu hồi quyền truy cập.
  \item Xóa đối tượng IAM.
\end{enumerate}

\subsubsection{Nhóm kịch bản quản lý EC2}
\subsubsection{Nhóm kịch bản quản lý Lambda}


\section{KIẾN TRÚC HỆ THỐNG DỰ KIẾN}

\subsection{Các thành phần và dịch vụ sử dụng}

\begin{table}[htbp]
\centering
\caption{Các thành phần và dịch vụ đề xuất trong kiến trúc hệ thống}
\label{tab:kientruc}
\bangfont
\begin{tabularx}{\linewidth}{@{}p{4.5cm}Y@{}}
\toprule
\textbf{Thành phần / Dịch vụ} & \textbf{Vai trò dự kiến}\\
\midrule
React & Xây dựng giao diện ứng dụng web (Frontend).\\
Amazon S3 \& Amazon CloudFront & Lưu trữ static file và phân phối giao diện web đến người dùng.\\
Amazon Cognito & Quản lý xác thực và đăng nhập người dùng ứng dụng web.\\
Amazon API Gateway & Tiếp nhận và định tuyến các yêu cầu HTTP/HTTPS từ Frontend đến Backend.\\
AWS Lambda & Chạy phần xử lý trung tâm (Backend serverless) của ứng dụng.\\
Python / Boto3 & Ngôn ngữ và thư viện AWS SDK được triển khai trong Lambda Backend để gọi các AWS API tương ứng.\\
LLM & Mô hình ngôn ngữ lớn phân tích ngữ cảnh yêu cầu của người dùng và hỗ trợ cơ chế \textit{function calling} (chưa chốt nhà cung cấp hoặc mô hình cụ thể).\\
Amazon DynamoDB & Cơ sở dữ liệu NoSQL lưu trữ lịch sử hội thoại, các tham số và trạng thái tác vụ.\\
Amazon CloudWatch & Ghi log hệ thống và theo dõi hoạt động vận hành của các dịch vụ backend.\\
AWS IAM & Thiết lập, cấp phát và kiểm soát quyền truy cập giữa các thành phần trong hệ thống.\\
\bottomrule
\end{tabularx}
\end{table}

\subsection{UML kiến trúc hệ thống / UML luồng}


\section{PHÂN CÔNG VÀ KẾ HOẠCH TRIỂN KHAI}

\subsection{Kế hoạch và phân công}
Kế hoạch triển khai được chia làm 4 sprint tương ứng với 4 tuần thực hiện.

\begin{table}[htbp]
\centering
\caption{Kế hoạch mục tiêu và đầu ra theo từng Sprint}
\label{tab:sprint_overview}
\bangfont
\begin{tabularx}{\linewidth}{@{}p{2.5cm}p{5cm}Y@{}}
\toprule
\textbf{Sprint} & \textbf{Mục tiêu chính} & \textbf{Đầu ra dự kiến}\\
\midrule
Sprint 1 & Thiết lập quy trình làm việc và kiểm thử PoC độc lập. & -- Thống nhất quy tắc làm việc.\\
& & -- Tạo AWS Account chung và danh tính IAM User cho từng thành viên.\\
& & -- Viết và kiểm thử độc lập các hàm Python/Boto3 cơ bản cho dịch vụ phụ trách.\\
\midrule
Sprint 2 & Dựng hạ tầng nền tảng và tích hợp Backend. & -- Giao diện React cơ bản trên S3/CloudFront.\\
& & -- Luồng đăng nhập người dùng qua Cognito.\\
& & -- API Gateway và Lambda Backend trung tâm tiếp nhận request.\\
& & -- Tích hợp các module Boto3 của 04 thành viên vào Backend.\\
\midrule
Sprint 3 & Tích hợp AI Chatbot và cơ chế Function Calling. & -- Tích hợp LLM (OpenAI/Anthropic) kèm prompt và cấu hình function calling.\\
& & -- Hoàn thiện luồng: Người dùng chat $\rightarrow$ Backend gọi LLM $\rightarrow$ LLM đề xuất tool/tham số $\rightarrow$ Backend xác thực và gọi Boto3 $\rightarrow$ AWS thực thi $\rightarrow$ Trả kết quả lên giao diện.\\
\midrule
Sprint 4 & Kiểm thử, tối ưu hóa và tổng kết dự án. & -- Kiểm thử các kịch bản thành công và các trường hợp lỗi/thiếu quyền.\\
& & -- Báo cáo tổng kết, slide thuyết trình hoàn thiện.\\
& & -- Dọn dẹp tài nguyên AWS thử nghiệm không còn cần thiết.\\
\bottomrule
\end{tabularx}
\end{table}

\begin{table}[htbp]
\centering
\caption{Bảng phân công trách nhiệm dự kiến cho từng thành viên}
\label{tab:phancong}
\bangfont
\begin{tabularx}{\linewidth}{@{}p{3cm}p{4.5cm}Y@{}}
\toprule
\textbf{Thành viên} & \textbf{Chuyên môn chính} & \textbf{Trách nhiệm phối hợp và thực hiện theo Sprint}\\
\midrule
Thành viên 1 & Phụ trách chính dịch vụ \textbf{S3} & -- \textbf{Sprint 1:} Viết và kiểm thử độc lập module Boto3 cho S3.\\
& & -- \textbf{Sprint 2:} Phối hợp dựng giao diện React và tích hợp module S3 vào Backend.\\
& & -- \textbf{Sprint 3:} Phối hợp định nghĩa function calling cho S3.\\
& & -- \textbf{Sprint 4:} Kiểm thử hệ thống, tối ưu module S3, làm báo cáo và slide.\\
\midrule
Thành viên 2 & Phụ trách chính dịch vụ \textbf{EC2} & -- \textbf{Sprint 1:} Thống nhất kịch bản EC2, viết và kiểm thử độc lập module Boto3 cho EC2.\\
& & -- \textbf{Sprint 2:} Phối hợp cấu hình Cognito/API Gateway và tích hợp module EC2.\\
& & -- \textbf{Sprint 3:} Phối hợp định nghĩa function calling cho EC2 và ghép luồng LLM.\\
& & -- \textbf{Sprint 4:} Kiểm thử hệ thống, tối ưu module EC2, dọn dẹp tài nguyên AWS.\\
\midrule
Thành viên 3 & Phụ trách chính dịch vụ \textbf{Lambda} & -- \textbf{Sprint 1:} Thống nhất kịch bản Lambda, viết và kiểm thử độc lập module Boto3 cho Lambda.\\
& & -- \textbf{Sprint 2:} Phụ trách chính dựng Lambda Backend trung tâm và tích hợp các module.\\
& & -- \textbf{Sprint 3:} Phối hợp tích hợp API LLM, xử lý phản hồi function calling tại Backend.\\
& & -- \textbf{Sprint 4:} Kiểm thử phân quyền, tối ưu hệ thống và chuẩn bị slide.\\
\midrule
Thành viên 4 & Phụ trách chính dịch vụ \textbf{IAM} & -- \textbf{Sprint 1:} Viết và kiểm thử độc lập module Boto3 cho IAM; hỗ trợ tạo IAM User cho nhóm.\\
& & -- \textbf{Sprint 2:} Cấu hình IAM Role cho Backend và triển khai S3/CloudFront cho Frontend.\\
& & -- \textbf{Sprint 3:} Hoàn thiện prompt và định nghĩa function calling cho IAM.\\
& & -- \textbf{Sprint 4:} Kiểm thử các trường hợp người dùng thiếu quyền, tổng hợp báo cáo.\\
\bottomrule
\end{tabularx}
\end{table}

\end{document}